% title:          Radius of curvature via centripetal acceleration
% area:           mechanics, differential-geometry
% methods:
% difficulty:
% difficulty-ai:  easy
% status:
% review:         none
% --- history: all optional, leave EMPTY if unknown ---
% origin:         folklore
% origin-date:
% origin-number:
% also-in:
% taken-from:
% references:     Standard calculus formula R = (1+y'^2)^{3/2}/|y''| for the curvature radius of a graph - https://en.wikipedia.org/wiki/Radius_of_curvature; the kinematic method (normal acceleration v^2/R along any curved trajectory) is a standard olympiad-physics idea, e.g. J. Kalda, Problems on Kinematics (IPhO study guide, version 2017), fact 13 and problem 24 - https://web.archive.org/web/2019/https://www.ioc.ee/~kalda/ipho/kin_ENG.pdf; the same derivation with constant dx/dt is discussed in Physics Stack Exchange Q715313 (2022) - https://physics.stackexchange.com/questions/715313; related constant-speed exercise Irodov, Problems in General Physics, Problem 1.42 (as quoted in Physics SE Q552788, 2020 - https://physics.stackexchange.com/questions/552788)
% history:        ai
\begin{problem}
Consider a particle moving along a curve $y=f(x)$. It is forced to move at a unit speed along the x direction. Find the expression of radius of curvature $R$, using the centripetal acceleration $\vec{a}_{\perp}$. 
\[ |\vec{a}_{\perp}|=\frac{|\vec{v}|^2}{R}\]
\end{problem}
